\section{Conclusion}
\label{sec:conclusion}

This thesis presents a comprehensive study on time series forecasting through the development of \textbf{MR-CDM}, a novel framework that addresses key limitations of existing methods. The research makes several significant contributions: (1) a delay embedding technique that converts variable-length time series into structured 2D image representations while preserving temporal dependencies, thereby enabling the use of spatial inductive biases from computer vision; (2) a hierarchical trend decomposition module that explicitly captures multi-scale temporal patterns—including short-term fluctuations, seasonal cycles, and long-term trends; and (3) a hierarchical conditional diffusion model that performs denoising generation under multi-scale guidance, reducing stochasticity through coarse-to-fine constraints.

% These components are supported by theoretical analyses—covering approximation guarantees, time complexity, and space complexity—as well as extensive empirical validation on real-world datasets, which demonstrate statistical consistency, computational feasibility, and consistent improvements of 15–20\% in prediction accuracy over state-of-the-art baselines.

% Looking ahead, several promising directions for future work emerge from this research. These include developing adaptive hyperparameter optimization strategies, incorporating exogenous variables (e.g., external covariates or contextual signals), enhancing uncertainty quantification for reliable probabilistic forecasting, improving computational efficiency via model compression or accelerated diffusion sampling, and exploring cross-domain generalization in high-impact applications such as healthcare monitoring, financial forecasting, and climate modeling. Overall, MR-CDM advances the field of time series forecasting by unifying image-based representation learning, multi-resolution temporal analysis, and conditional generative modeling. It not only achieves state-of-the-art performance on standard benchmarks but also offers a new paradigm for tackling fundamental challenges in temporal data modeling, with substantial potential for both academic innovation and industrial deployment.
    